% ===================================================================== CHƯƠNG 2
\chapter{Cơ sở lý thuyết}\label{ch:co-so}

Chương này không tổng quan toàn bộ lĩnh vực. Nó chỉ trình bày những kiến thức làm căn cứ trực tiếp cho các lựa chọn thiết kế ở Chương~\ref{ch:thiet-ke}. Mỗi mục kết thúc bằng một hệ quả cho thiết kế. Nhờ vậy, người đọc thấy mỗi quyết định kỹ thuật bắt nguồn từ đâu, thay vì chỉ là sở thích của người làm.

\section{Tẩy xanh: định nghĩa và hệ quả cho thiết kế}

Delmas và Burbano \cite{delmas2011} định nghĩa tẩy xanh là sự kết hợp giữa thành tích môi trường kém và truyền thông tích cực về thành tích đó. Lyon và Montgomery \cite{lyon2015} cho rằng thuật ngữ này đã thành một “chiếc ô” bao trùm nhiều hành vi khác nhau, từ chọn lọc thông tin tới tuyên bố mơ hồ. Seele và Gatti \cite{seele2017} chuyển trọng tâm sang người quan sát: tẩy xanh tồn tại khi có \emph{cáo buộc} từ bên ngoài, và cần phân biệt với cáo buộc oan. Báo cáo của ba cơ quan giám sát châu Âu \cite{esas2024} nhấn mạnh rằng tẩy xanh có thể vô ý.

Ở mức thực hành, nhóm “bảy tội” của TerraChoice \cite{terrachoice2009} cho thấy những dạng có thể nhận ra từ chính văn bản: tuyên bố không có bằng chứng, tuyên bố mơ hồ, tuyên bố không liên quan, và đánh đổi bị che giấu. Hai dạng đầu là những dạng mà một công cụ đọc văn bản có thể chỉ ra một cách khách quan: câu thiếu số liệu, thiếu kỳ, thiếu phạm vi, hoặc dùng từ ngữ chung chung.

\begin{hopy}{Hệ quả cho thiết kế}
Không gán nhãn tẩy xanh cho văn bản hay doanh nghiệp. Thay vào đó, kiểm từng tuyên bố theo những gì đo được: có bằng chứng không, bằng chứng có so được không, câu thiếu thuộc tính nào, có dùng từ mơ hồ không. Kết luận cuối thuộc về người soát xét.
\end{hopy}

\section{Chuẩn mực đảm bảo và chuẩn mực công bố}

Dịch vụ đảm bảo thông tin phi tài chính được thực hiện theo ISAE 3000 (sửa đổi) \cite{isae3000}. Chuẩn mực mới ISSA 5000 \cite{issa5000}, ban hành tháng 11/2024, áp dụng cho dịch vụ đảm bảo thông tin bền vững từ kỳ bắt đầu ngày 15/12/2026. Cả hai xoay quanh một khái niệm: người thực hiện phải thu thập \emph{bằng chứng đủ và thích hợp} cho thông tin được đảm bảo, và lưu hồ sơ để người khác tái dựng được.

Về phía công bố, GHG Protocol \cite{ghgprotocol} quy định cách báo cáo phát thải theo phạm vi 1, 2, 3, theo ranh giới tổ chức (công ty mẹ hay hợp nhất) và theo năm gốc. IFRS S2 \cite{ifrss2} và GRI 305 yêu cầu nêu rõ chỉ số, kỳ, phạm vi, phương pháp. Theo các chuẩn mực này, hai con số chỉ so được với nhau khi chúng cùng chỉ số, cùng kỳ, cùng phạm vi, cùng ranh giới và cùng cơ sở đo. Ví dụ: không thể đem phát thải phạm vi 1 của một nhà máy so với phát thải phạm vi 1 và 2 của cả tập đoàn.

\begin{hopy}{Hệ quả cho thiết kế}
Năm thuộc tính bắt buộc của một tuyên bố (chỉ số, giá trị, kỳ, năm gốc, phạm vi) lấy từ những gì chuẩn mực đòi hỏi để một con số có thể kiểm được. Hệ thống phải có một \emph{cổng so sánh}: chỉ đem hai con số ra so khi chúng thực sự cùng đại lượng. Hệ thống tạo ra hồ sơ làm việc có trang nguồn và phép tính, nhưng không tự nhận là đáp ứng ISAE 3000 hay ISSA 5000.
\end{hopy}

\section{Khung pháp lý Việt Nam liên quan}

Bảng~\ref{tab:phap-ly} tóm tắt các văn bản mà hệ thống dùng làm nguồn tham chiếu pháp lý. Mỗi văn bản được lưu cùng ngày ký, ngày hiệu lực và chuỗi văn bản sửa đổi. Lý do là một văn bản chỉ áp dụng được cho một tuyên bố khi nó còn hiệu lực trong kỳ của tuyên bố đó, đúng ngành và đúng vấn đề.

\begin{table}[htbp]
\centering
\caption{Các văn bản pháp lý chính trong sổ đăng ký pháp lý của hệ thống}\label{tab:phap-ly}
\bangchu
\rowcolors{2}{rowbg}{white}
\begin{tabularx}{\linewidth}{L{4.3cm}Y}
\hdr\thd{Văn bản} & \thd{Nội dung liên quan tới kiểm chứng tuyên bố}\\
Luật Bảo vệ môi trường 72/2020/QH14 & Khung chung; trái phiếu xanh (Điều 150); sản phẩm “thân thiện môi trường” phải được chứng nhận.\\
Nghị định 08/2022/NĐ-CP, sửa bởi 05/2025/NĐ-CP & Quy định chi tiết Luật Bảo vệ môi trường.\\
Nghị định 06/2022/NĐ-CP, sửa bởi 119/2025/NĐ-CP & Giảm nhẹ phát thải khí nhà kính, kiểm kê và thẩm định; bảo vệ tầng ô-dôn.\\
Quyết định 13/2024/QĐ-TTg & Danh mục lĩnh vực, cơ sở phát thải phải kiểm kê khí nhà kính.\\
Quyết định 21/2025/QĐ-TTg & Tiêu chí môi trường và xác nhận dự án thuộc danh mục phân loại xanh.\\
Thông tư 96/2020/TT-BTC & Công bố thông tin trên thị trường chứng khoán; chỉ tiêu môi trường trong báo cáo thường niên.\\
Thông tư 17/2022/TT-NHNN & Quản lý rủi ro môi trường trong hoạt động cấp tín dụng.\\
Nghị định 48/2026/NĐ-CP, 83/2026/NĐ-CP & Sửa đổi các nghị định trên; Nghị định 83/2026 chỉ sửa thủ tục về chất được kiểm soát (ô-dôn, HFC).\\
\end{tabularx}
\end{table}

Trường hợp Nghị định 83/2026 là một bài học thật. Trước ngày 06/10/2026, văn bản này bị gắn nhầm chủ đề “phát thải” nên được trích kèm 55 tuyên bố phát thải, trong khi nó chỉ sửa thủ tục về chất ô-dôn và HFC. Lỗi được sửa bằng cách gắn mỗi văn bản sửa đổi đúng với vấn đề nó sửa.

\begin{hopy}{Hệ quả cho thiết kế}
Áp dụng pháp lý phải kiểm \emph{hiệu lực áp dụng} (ngày, ngành, chủ thể, loại tuyên bố), không chỉ kiểm “cùng chủ đề”. Quy tắc pháp lý do người có chuyên môn viết và duyệt, có phiên bản và có kiểm thử. Mô hình không được tự đọc luật rồi sinh quy tắc trong lúc chạy.
\end{hopy}

\section{Kiểm chứng tuyên bố tự động}

Bài toán kiểm chứng tuyên bố trong xử lý ngôn ngữ tự nhiên thường gồm ba bước: tìm bằng chứng, đánh giá quan hệ giữa tuyên bố và bằng chứng, rồi đưa ra kết luận \cite{guo2022survey}. Bộ dữ liệu FEVER \cite{thorne2018} dùng ba nhãn: được ủng hộ, bị bác bỏ và \emph{không đủ thông tin}. CLIMATE-FEVER \cite{diggelmann2020} áp dụng cách làm này cho tuyên bố khí hậu. Trong lĩnh vực môi trường doanh nghiệp, ClimateBERT \cite{webersinke2022} là mô hình ngôn ngữ được huấn luyện thêm trên văn bản khí hậu. Bộ dữ liệu Environmental Claim Detection \cite{stammbach2023} tách bài toán \emph{nhận diện câu nào là tuyên bố môi trường} thành một tác vụ riêng. Khảo sát của Calamai và cộng sự \cite{calamai2025} nhận xét rằng các nghiên cứu hiện có chủ yếu giải các tác vụ trung gian như vậy, và khuyến nghị thiết kế phân rã, dễ giải thích, có người giám sát.

Với số liệu tài chính và môi trường, khó khăn lớn nằm ở suy luận số. Các bộ FinQA \cite{chen2021finqa} và TAT-QA \cite{zhu2021tatqa} cho thấy mô hình ngôn ngữ thường sai ở các bước tính toán và đọc bảng. FinDVer \cite{zhao2024findver}, bộ kiểm chứng tuyên bố trên tài liệu tài chính dài, cho thấy các mô hình lớn vẫn còn cách xa chuyên gia khi phải kiểm tuyên bố có số liệu.

\begin{hopy}{Hệ quả cho thiết kế}
Tách hệ thống thành các tác vụ con đo được: nhận diện tuyên bố, đọc thuộc tính, tìm bằng chứng, kiểm chứng. Giữ nhãn “chưa đủ bằng chứng” như một kết luận hợp lệ. Giao việc tính toán cho mã tất định, không giao cho mô hình ngôn ngữ.
\end{hopy}

\section{Mô hình ngôn ngữ lớn: năng lực đọc và các rủi ro}

Kỹ thuật kết hợp truy xuất với sinh văn bản \cite{lewis2020rag} giúp mô hình trả lời dựa trên tài liệu. Tuy vậy, có bốn rủi ro được ghi nhận rõ trong tài liệu khoa học và đều liên quan trực tiếp tới một công cụ dùng trong kiểm toán:

\begin{itemize}
\item \textbf{Bịa và tính sai.} Mô hình có thể đưa ra con số không có trong nguồn, hoặc tính sai phần trăm và đổi sai đơn vị (mục trước).
\item \textbf{Lời giải thích không trung thực.} Turpin và cộng sự \cite{turpin2023} cho thấy lời giải thích từng bước của mô hình có thể không phản ánh lý do thật dẫn tới câu trả lời.
\item \textbf{Mô hình làm người chấm có thiên lệch.} Mô hình dùng làm người chấm có thể thiên vị vị trí, độ dài, hoặc văn phong của chính họ mô hình đó \cite{zheng2023judge,wang2024fair}.
\item \textbf{Tiêm lệnh gián tiếp.} Một câu lệnh giấu trong tài liệu có thể điều khiển mô hình đang đọc tài liệu đó \cite{greshake2023}. Các nghiên cứu về thứ bậc chỉ dẫn \cite{wallace2024} cho thấy cần phân biệt rõ đâu là chỉ dẫn của hệ thống và đâu là dữ liệu.
\end{itemize}

\begin{hopy}{Hệ quả cho thiết kế}
Mô hình chỉ \emph{đọc}: chuyển văn bản thành biểu diễn có cấu trúc, gợi ý đoạn nào đáng xem. Mô hình không đứng một mình mà ra được kết luận dứt khoát. Tài liệu người dùng đưa vào luôn là \emph{dữ liệu}, không bao giờ là chỉ dẫn. Khi dùng mô hình để chấm, kết quả chấm phải được ghi rõ là chấm bằng máy và có thể thiên lệch.
\end{hopy}

\section{Dự đoán có từ chối và đo trên mẫu nhỏ}

Lý thuyết dự đoán có từ chối \cite{chow1970,geifman2017} chỉ ra một sự đánh đổi: hệ thống càng chỉ kết luận ở những ca chắc chắn thì càng ít sai trên phần đã kết luận, nhưng càng nhường nhiều ca cho người. Đây chính là bài toán của người soát xét. Một kết luận sai về doanh nghiệp đắt hơn nhiều so với một ca phải đọc thêm. Vì vậy mục tiêu hợp lý là \emph{tối đa hoá phần kết luận được, với giới hạn trên số kết luận sai nguy hiểm}, chứ không phải tối đa hoá một độ chính xác chung. Hiệu chỉnh xác suất \cite{guo2017calibration} là bước cần có nếu muốn từ chối dựa trên độ tin cậy của mô hình.

Khi mẫu nhỏ, cách báo con số quan trọng không kém con số. Báo cáo dùng các công cụ sau.
\begin{itemize}
\item \textbf{Khoảng tin cậy Wilson} \cite{wilson1927} cho tỷ lệ. Ví dụ: 0 lỗi trên 87 ca vẫn cho phép tỷ lệ lỗi thật lên tới khoảng 3,4\%.
\item \textbf{Bootstrap theo công ty}, vì các câu trong cùng một báo cáo không độc lập với nhau.
\item \textbf{Hệ số đồng thuận κ} \cite{cohen1960}, diễn giải theo thang của Landis và Koch \cite{landis1977}, cùng biến thể có trọng số bậc hai cho thang thứ bậc.
\end{itemize}

\begin{hopy}{Hệ quả cho thiết kế}
Báo cáo tách ba chỉ số: độ phủ, độ chính xác trên phần dám kết luận, và số kết luận sai nguy hiểm theo từng loại. Mỗi con số đi kèm cỡ mẫu, khoảng tin cậy và loại căn cứ đối chiếu.
\end{hopy}

\section{Học từ nhãn do máy tạo}

Gilardi và cộng sự \cite{gilardi2023} cho thấy mô hình ngôn ngữ lớn có thể gán nhãn văn bản với chất lượng ngang hoặc hơn người gán thuê ở nhiều tác vụ, với chi phí thấp hơn nhiều. Điều này mở ra cách làm \emph{chưng cất}: dùng mô hình lớn (hoặc một quy trình có máy hỗ trợ) để tạo nhãn, rồi huấn luyện một mô hình nhỏ chạy tại máy. Mô hình mã hoá đa ngôn ngữ như Multilingual E5 \cite{wang2024e5} đã hiểu nghĩa tiếng Việt từ trước, nên chỉ cần vài trăm tới vài nghìn câu có nhãn để học ranh giới giữa các lớp.

Cách làm này có ba rủi ro. Mô hình học trò thừa hưởng sai lệch có hệ thống của nguồn nhãn. Hai lượt gán bằng cùng một công cụ có thể sai giống nhau, nên độ đồng thuận cao giữa hai lượt chỉ cho biết nhãn \emph{nhất quán}, không cho biết nhãn \emph{đúng}. Và người rà soát nhãn máy dễ bị neo vào nhãn có sẵn.

\begin{hopy}{Hệ quả cho thiết kế}
Nhãn có máy hỗ trợ được dùng để huấn luyện và để dò lỗi lớn nhanh. Độ chính xác cuối cùng chỉ được đo trên nhãn chuẩn do người gán độc lập, không xem nhãn máy. Bộ nhãn chuẩn không bao giờ dùng để huấn luyện.
\end{hopy}

\section{Vị trí của công nghệ lượng tử}\label{sec:luong-tu}

Cuộc thi có nhóm đề tài lượng tử, và thuật toán tối ưu xấp xỉ lượng tử (QAOA) \cite{farhi2014} cùng các bài toán dạng tối ưu nhị phân bậc hai không ràng buộc (QUBO) thường được nhắc tới trong ứng dụng tài chính \cite{herman2023}. Bài toán cốt lõi của báo cáo là \emph{đọc hiểu và kiểm chứng văn bản}, không phải bài toán tối ưu tổ hợp lớn. Ở quy mô một báo cáo vài trăm trang, chưa có bằng chứng cho thấy lợi thế lượng tử. Ứng dụng duy nhất có thể hợp lý là chọn mẫu tuyên bố để soát theo ngân sách thời gian. Ứng dụng này mới được hình thức hoá trên giấy, chưa có dòng mã nào.

\begin{hopy}{Hệ quả cho thiết kế}
Không thêm thành phần lượng tử chỉ để có. Thể lệ cuộc thi cũng nêu rõ việc dùng công nghệ phức tạp hay mới không phải là tiêu chí tự thân. Bài toán chọn mẫu soát theo QUBO chỉ được xem xét khi chứng minh được giá trị so với thuật toán cổ điển trên dữ liệu thật.
\end{hopy}

\FloatBarrier
\begin{hopy}{Tóm tắt chương}
\begin{itemize}
\item Tẩy xanh là một cáo buộc, không phải thuộc tính của văn bản. Hệ thống kiểm từng tuyên bố theo những gì đo được.
\item Chuẩn mực đảm bảo và công bố quyết định năm thuộc tính bắt buộc và quy tắc “chỉ so cái so được”.
\item Pháp lý phải kiểm hiệu lực áp dụng; quy tắc do người viết và duyệt.
\item Mô hình ngôn ngữ đọc tốt nhưng có thể bịa, tính sai, giải thích không trung thực và bị tiêm lệnh. Vì vậy mô hình đọc, quy tắc quyết định, người duyệt.
\item Dự đoán có từ chối là khung đúng cho bài toán; nhãn máy để dạy, nhãn người để đo; lượng tử không có chỗ trong lõi ở quy mô hiện tại.
\end{itemize}
\end{hopy}
